% © 2026 Ing. David Jaroš — CC BY-NC-ND 4.0
%
% This work is licensed under a Creative Commons Attribution-NonCommercial-NoDerivatives
% 4.0 International License (CC BY-NC-ND 4.0).
%
% Kac-Moody Level k from UBT Action S[Θ] — Approaches H1, H2, H3
% Track: CORE — α derivation — v67 baseline
% Date: 2026-03-08

\ifdefined\INCLUDEMODE
  % Being included in another document — skip preamble
\else
  \documentclass[11pt,a4paper]{article}
  \usepackage{amsmath,amssymb,amsthm}
  \usepackage{hyperref}
  \usepackage{booktabs}

  \newtheorem{theorem}{Theorem}
  \newtheorem{lemma}{Lemma}
  \newtheorem{proposition}{Proposition}
  \newtheorem{conjecture}{Conjecture}
  \theoremstyle{definition}
  \newtheorem{gap}{Gap}
  \theoremstyle{remark}
  \newtheorem{remark}{Remark}
  \newtheorem{observation}{Observation}

  \title{Kac-Moody Level $k$ from UBT Action $\mathcal{S}[\Theta]$\\
    \large Approaches H1 (WZW Normalization), H2 (Absence of CS Term),\\
    H3 (Dynkin Index): v67 Baseline}
  \author{Ing.\ David Jaroš}
  \date{2026-03-08 (v67)}

  \begin{document}
  \maketitle
\fi

%--------------------------------------------------------------------
\section{Problem Statement and Context}
\label{sec:intro}
%--------------------------------------------------------------------

\noindent\textbf{No circular reasoning.}
$\alpha^{-1} = 137$ and $n^* = 137$ are \emph{never} used as inputs
anywhere in this document.
The inputs are: the UBT kinetic action
$\mathcal{S}[\Theta] = \int \mathrm{Sc}[(D_\mu\Theta)^\dagger(D^\mu\Theta)]\,\mathrm{d}^4x\,\mathrm{d}\psi$
with $\psi \in [0, 2\pi)$, the algebraic structure $\mathbb{C}\otimes\mathbb{H}$,
and $N_{\mathrm{eff}} = 12$ (proved independently of $\alpha$).

\subsection{Why is the Kac-Moody level $k$ needed?}

Approach~G3 in \texttt{b\_base\_g\_approaches.tex} (v65) established a
critical observation:

\begin{observation}[G3 level-$k=1$ match, from v65]
  \label{obs:G3_k1_recap}
  If the $\psi$-circle of UBT carries a level-$k=1$
  $\widehat{\mathrm{SU}(2)}$ Kac-Moody symmetry, then
  $c_{\widehat{\mathrm{SU}(2)},1} = 1$ and the formula
  $B_{\mathrm{base}} = c \cdot N_{\mathrm{eff}}^{3/2}$ reproduces
  $B_{\mathrm{base}} = 41.57$ exactly with no free parameters.
  The selection of $k = 1$ was motivated by free-boson / minimal-coupling
  arguments, but was not uniquely derived from UBT axioms.
  Status: \textbf{[PARTIAL]}.
\end{observation}

The central charge formula for $\widehat{\mathrm{SU}(2)}_k$ is:
\begin{equation}
  c_{\mathrm{SU}(2),k}
  \;=\; \frac{3k}{k+2},
  \label{eq:c_su2_k}
\end{equation}
which gives $c = 1$ for $k = 1$ and $c = 18/7 \approx 2.571$ for
$k = 12$, with all other levels lying between these extremes or above.
The unique level giving $c = 1$ is $k = 1$.

\textbf{This document has one goal:}
Compute the level $k$ of the Kac-Moody algebra $\widehat{\mathrm{SU}(2)}$
on the $\psi$-circle directly from $\mathcal{S}[\Theta]$, without using
$\alpha^{-1} = 137$ or $n^* = 137$ as inputs.

Three approaches are investigated:
\begin{itemize}
  \item \textbf{H1}: WZW normalization — compare the UBT kinetic term
    with the standard WZW action to read off $k$ from $\langle|\Theta|^2\rangle_\mathrm{vac}$.
  \item \textbf{H2}: Absence of Chern-Simons term — parity argument that
    $k_{\mathrm{CS}} = 0$, and the free-field / minimal-coupling argument
    that the kinetic term alone sets $k = 1$.
  \item \textbf{H3}: Dynkin index — compute $k$ from the number and
    representation content of the modes contributing to the
    $\widehat{\mathrm{SU}(2)}$ current algebra.
\end{itemize}

%--------------------------------------------------------------------
\section{Setup: SU(2) Structure in \texorpdfstring{$\mathbb{C}\otimes\mathbb{H}$}
{C⊗H} and the WZW Framework}
\label{sec:setup}
%--------------------------------------------------------------------

\subsection{SU(2) subgroup of \texorpdfstring{$\mathbb{C}\otimes\mathbb{H}$}{C⊗H}}

The UBT field $\Theta \in \mathbb{C}\otimes\mathbb{H} \cong \mathrm{Mat}(2,\mathbb{C})$.
We decompose:
\begin{equation}
  \Theta = \Theta_0 \cdot \mathbf{1} + \Theta_1 \cdot \mathbf{i}
         + \Theta_2 \cdot \mathbf{j} + \Theta_3 \cdot \mathbf{k},
  \quad \Theta_\mu \in \mathbb{C},
  \label{eq:Theta_decomp}
\end{equation}
with $\mathrm{Sc}[\Theta^\dagger\Theta] = |\Theta_0|^2 + |\Theta_1|^2
+ |\Theta_2|^2 + |\Theta_3|^2$.

The \emph{compact real form} $\mathrm{SU}(2) \subset \mathrm{Mat}(2,\mathbb{C})$
consists of elements with $\Theta_\mu \in \mathbb{R}$ and
$\Theta_0^2 + \Theta_1^2 + \Theta_2^2 + \Theta_3^2 = 1$.
These are unit quaternions parameterising $S^3$.

For a general $\Theta$ with vacuum amplitude $r = \sqrt{\mathrm{Sc}[\Theta^\dagger\Theta]}$,
we write:
\begin{equation}
  \Theta = r \cdot g, \quad g = \frac{\Theta}{|\Theta|} \in \mathrm{SU}(2),
  \quad r \in \mathbb{R}_{>0}.
  \label{eq:Theta_polar}
\end{equation}
This polar decomposition is well-defined when $\Theta$ is in the quaternionic
(not fully complex) subspace.

\subsection{WZW action on the \texorpdfstring{$\psi$}{psi}-circle}

The standard WZW model on a circle $S^1$ for group $\mathrm{SU}(2)$ at level $k$
has the kinetic term (before the topological WZ term, which vanishes on $S^1$):
\begin{equation}
  S_{\mathrm{WZW}}^{S^1}[g]
  \;=\; -\frac{k}{4\pi} \int_0^{2\pi} \mathrm{d}\psi\;
    \mathrm{Tr}_{\mathrm{fund}}\!\bigl[(g^{-1}\partial_\psi g)^2\bigr],
  \label{eq:WZW_S1}
\end{equation}
where $\mathrm{Tr}_{\mathrm{fund}}$ is the trace in the fundamental (spin-$\frac{1}{2}$)
representation, normalised so that
$\mathrm{Tr}_{\mathrm{fund}}(T^a T^b) = \tfrac{1}{2}\delta^{ab}$ for
$T^a = \sigma^a/2$.
The Lie algebra element $J_\psi := g^{-1}\partial_\psi g \in \mathfrak{su}(2)$
is anti-Hermitian; hence $-\mathrm{Tr}_{\mathrm{fund}}[J_\psi^2] > 0$.

%--------------------------------------------------------------------
\section{Approach H1: WZW Normalization}
\label{sec:H1}
%--------------------------------------------------------------------

\textbf{Status: [BLOCKED BY CIRCULARITY]} — see Observation~\ref{obs:H1_circularity}.

\subsection{Derivation of the kinetic term for the SU(2) sector}

Substituting $\Theta = r \cdot g$ into the UBT kinetic action and
restricting to $\psi$-derivatives (the 4D spacetime integral is treated
as an overall volume factor):
\begin{equation}
  \mathcal{S}_{\mathrm{kin}}^\psi[\Theta]
  \;=\; \int_0^{2\pi} \mathrm{d}\psi\;
    \mathrm{Sc}\bigl[(\partial_\psi\Theta)^\dagger(\partial_\psi\Theta)\bigr].
  \label{eq:S_kin_psi}
\end{equation}
Expanding with $\Theta = r g$:
\begin{equation}
  \partial_\psi\Theta = (\partial_\psi r)\,g + r\,\partial_\psi g,
  \label{eq:dTheta_expand}
\end{equation}
\begin{equation}
  (\partial_\psi\Theta)^\dagger(\partial_\psi\Theta)
  = (\partial_\psi r)^2\, g^\dagger g
  + r^2\, (\partial_\psi g)^\dagger(\partial_\psi g)
  + r(\partial_\psi r)\bigl[(\partial_\psi g)^\dagger g + g^\dagger\partial_\psi g\bigr].
  \label{eq:dTheta_sq}
\end{equation}
The cross-term vanishes when taking the scalar part, because
$g^\dagger \partial_\psi g \in \mathfrak{su}(2)$ is anti-Hermitian and
$\mathrm{Sc}[A] = 0$ for any anti-Hermitian $A$.
Using $\mathrm{Sc}[g^\dagger g] = \mathrm{Sc}[\mathbf{1}] = 1$:
\begin{equation}
  \mathcal{S}_{\mathrm{kin}}^\psi[\Theta]
  = \int_0^{2\pi} \mathrm{d}\psi\,\bigl[(\partial_\psi r)^2
  + r^2\,\mathrm{Sc}\bigl[(\partial_\psi g)^\dagger(\partial_\psi g)\bigr]\bigr].
  \label{eq:S_kin_split}
\end{equation}

\subsection{Relating \texorpdfstring{$\mathrm{Sc}$}{Sc} to the Killing norm}
\label{subsec:Sc_to_Tr}

For $A \in \mathrm{Mat}(2,\mathbb{C})$, the scalar part satisfies
$\mathrm{Sc}[A] = \tfrac{1}{2}\mathrm{Tr}(A)$ (for the $2\times 2$ matrix trace).
Therefore:
\begin{equation}
  \mathrm{Sc}\bigl[(\partial_\psi g)^\dagger(\partial_\psi g)\bigr]
  = \tfrac{1}{2}\,\mathrm{Tr}_{\mathrm{fund}}\bigl[(\partial_\psi g)^\dagger(\partial_\psi g)\bigr].
  \label{eq:Sc_Tr}
\end{equation}
For $g \in \mathrm{SU}(2)$, we have $J_\psi := g^{-1}\partial_\psi g \in \mathfrak{su}(2)$,
so $J_\psi^\dagger = -J_\psi$.
Writing $\partial_\psi g = g J_\psi$:
\begin{equation}
  (\partial_\psi g)^\dagger(\partial_\psi g)
  = J_\psi^\dagger g^\dagger g J_\psi
  = J_\psi^\dagger J_\psi = -J_\psi^2.
  \label{eq:dg_sq}
\end{equation}
Hence:
\begin{equation}
  \mathrm{Sc}\bigl[(\partial_\psi g)^\dagger(\partial_\psi g)\bigr]
  = -\tfrac{1}{2}\,\mathrm{Tr}_{\mathrm{fund}}[J_\psi^2]
  = \tfrac{1}{2}\,\mathrm{Tr}_{\mathrm{fund}}[J_\psi^\dagger J_\psi]
  > 0.
  \label{eq:Sc_positive}
\end{equation}
Substituting into Eq.~\eqref{eq:S_kin_split}, the $g$-sector of the
$\psi$-circle kinetic action is:
\begin{equation}
  \mathcal{S}_{\mathrm{kin}}^{g,\psi}
  = r^2 \int_0^{2\pi} \mathrm{d}\psi\;
    \Bigl(-\tfrac{1}{2}\,\mathrm{Tr}_{\mathrm{fund}}[J_\psi^2]\Bigr).
  \label{eq:S_kin_g}
\end{equation}

\subsection{Comparison with the WZW action}

Comparing Eq.~\eqref{eq:S_kin_g} with the WZW action
\eqref{eq:WZW_S1}, both of the form $-C \int \mathrm{Tr}_{\mathrm{fund}}[J_\psi^2]\,\mathrm{d}\psi$:
\begin{equation}
  \underbrace{r^2 \times \tfrac{1}{2}}_{\text{UBT coefficient}}
  \;=\;
  \underbrace{\frac{k}{4\pi}}_{\text{WZW coefficient}}.
  \label{eq:coefficient_match}
\end{equation}
Solving for $k$:
\begin{equation}
  \boxed{k \;=\; 2\pi\,r^2 \;=\; 2\pi\,\langle\mathrm{Sc}[\Theta^\dagger\Theta]\rangle_{\mathrm{vac}}.}
  \label{eq:k_from_r2}
\end{equation}
\textit{Why this formula matters:} if we know the vacuum expectation value of
the UBT field norm, Eq.~\eqref{eq:k_from_r2} directly gives the
Kac-Moody level.

\subsection{Evaluating \texorpdfstring{$\langle|\Theta|^2\rangle_{\mathrm{vac}}$}{⟨|Θ|²⟩\_vac}}

For $k = 1$:
\begin{equation}
  r^2 = \langle\mathrm{Sc}[\Theta^\dagger\Theta]\rangle_{\mathrm{vac}}
  \stackrel{!}{=} \frac{1}{2\pi} \approx 0.1592.
  \label{eq:r2_for_k1}
\end{equation}
For $k \in \mathbb{Z}_{>0}$, the vacuum value $r^2 = k/(2\pi)$
must be determined from the UBT effective potential $V_{\mathrm{eff}}$.

The effective potential $V_{\mathrm{eff}}(n)$ (derived in
\texttt{step2\_vacuum\_polarization.tex}) has its minimum at
$n^* = 137$, and $|{\Theta_{\mathrm{vac}}}|^2$ scales with $n^*$.
In other words:
\begin{equation}
  \langle\mathrm{Sc}[\Theta^\dagger\Theta]\rangle_{\mathrm{vac}}
  = f(n^*),
  \label{eq:r2_depends_on_nstar}
\end{equation}
where $f$ is a function of the vacuum mode number $n^*$.
Since $n^* = 137 = \alpha^{-1}$ (which is the result of the derivation,
not an input), using $n^*$ to determine $r^2$ and hence $k$ would
introduce exactly the circular reasoning we aim to avoid.

\begin{observation}[H1 circularity blockage]
  \label{obs:H1_circularity}
  The WZW normalization formula $k = 2\pi r^2$ expresses the
  Kac-Moody level in terms of the vacuum field amplitude $r^2$.
  However, $r^2 = f(n^*)$, where $n^*$ is the vacuum winding number
  whose value ($n^* = 137$) is the quantity being derived.
  Any attempt to compute $k$ through $r^2$ from $V_{\mathrm{eff}}$
  is therefore \textbf{circular}.
  \textbf{[DEAD END via $V_{\mathrm{eff}}$]}.
\end{observation}

\subsection{Can canonical normalization rescue H1?}

\noindent\textit{Canonical normalization hypothesis:}
Suppose the UBT field is canonically normalised in natural units so that
$\mathrm{Sc}[\Theta^\dagger\Theta] = 1$ (unit biquaternion norm on $S^3$).
Then $r^2 = 1$ and:
\begin{equation}
  k \;=\; 2\pi \times 1 \;=\; 2\pi \;\approx\; 6.28.
  \label{eq:k_canonical}
\end{equation}
This is not a positive integer.
The Kac-Moody quantisation condition requires $k \in \mathbb{Z}_{>0}$
for $\mathrm{SU}(2)$ (or $k \in \tfrac{1}{2}\mathbb{Z}_{>0}$ for the spin
extension), so $k = 2\pi$ is unphysical.

\begin{observation}[H1 canonical normalization failure]
  \label{obs:H1_canonical}
  The canonical normalization $\mathrm{Sc}[\Theta^\dagger\Theta] = 1$
  gives $k = 2\pi \notin \mathbb{Z}$, which is not a valid
  Kac-Moody level.
  No algebraically natural choice of $r^2$ from $\mathbb{C}\otimes\mathbb{H}$
  (such as $1$, $1/2$, $1/4$, $1/(2\pi)$, etc.) is
  singled out by $\mathcal{S}[\Theta]$ alone.
  \textbf{[DEAD END via canonical normalization]}.
\end{observation}

\textbf{H1 result summary:}
\begin{itemize}
  \item Via $V_{\mathrm{eff}}$ minimum:
    blocked by circularity through $n^* = 137$.
    \textbf{[DEAD END (circular)]}.
  \item Via canonical normalization $r^2 = 1$:
    $k = 2\pi \notin \mathbb{Z}$.
    \textbf{[DEAD END (non-integer level)]}.
  \item No other natural choice of $r^2$ from $\mathcal{S}[\Theta]$
    gives $k = 1$.
\end{itemize}

%--------------------------------------------------------------------
\section{Approach H2: Absence of Chern-Simons Term}
\label{sec:H2}
%--------------------------------------------------------------------

\textbf{Status: [MOTIVATED CONJECTURE — k=1 motivated by absence of CS
and minimality principle; not uniquely proved.]}

\subsection{The role of the Chern-Simons term in fixing the WZW level}

In gauge theories on a manifold with boundary, or in topological field
theories on an odd-dimensional manifold, the Chern-Simons (CS) term
\begin{equation}
  \mathcal{S}_{\mathrm{CS}}
  = \frac{k_{\mathrm{CS}}}{4\pi}
    \int \mathrm{Tr}_{\mathrm{fund}}\!\Bigl(A \wedge \mathrm{d}A
      + \tfrac{2}{3} A \wedge A \wedge A\Bigr)
  \label{eq:CS_action}
\end{equation}
contributes an integer shift to the Kac-Moody level: $k \to k + k_{\mathrm{CS}}$.
Quantisation of the CS term requires $k_{\mathrm{CS}} \in \mathbb{Z}$ for
the path integral to be well-defined.
\textit{Why this matters:} if $\mathcal{S}[\Theta]$ contains a CS term with
$k_{\mathrm{CS}} \neq 0$, the effective Kac-Moody level is shifted
from its kinetic value, and $k = 1$ would require a delicate
cancellation between the kinetic and CS contributions.

\subsection{Parity argument: CS is absent from UBT}

\begin{proposition}[CS term absent from UBT]
  \label{prop:no_CS}
  The UBT action $\mathcal{S}[\Theta] = \mathcal{S}_{\mathrm{kin}} + \mathcal{S}_{\mathrm{pot}}$
  does not contain a Chern-Simons term.
  \begin{proof}[Argument]
    \begin{enumerate}
      \item \textbf{CS is P-odd.}
        The CS term $\mathrm{Tr}(A \wedge \mathrm{d}A + \ldots)$ is
        odd under parity $P$ (it is a pseudo-scalar in 3D, or a
        pseudo-action on $S^1$).
      \item \textbf{UBT has parity symmetry.}
        The UBT potential $V_{\mathrm{eff}}(n)$ (proved in
        \texttt{step2\_vacuum\_polarization.tex}) has a principal
        minimum at $n^* = 137$ and a mirror minimum at $n^{**} = 139$
        (the CP-conjugate sector).
        The two sectors are related by $P$ (or $CP$), so parity is a
        symmetry of the full UBT theory.
        \textit{Why this implies no CS: a P-odd term in the action
        would break this symmetry; since parity is preserved globally,
        no P-odd term can appear.}
      \item \textbf{CS on $S^1_\psi$ vanishes for trivial holonomy.}
        Even if a CS-type coupling were present in the 5D action,
        its $\psi$-circle integral reduces to a topological term that
        vanishes when the gauge holonomy around $S^1_\psi$ is trivial.
        The Hosotani mechanism (proved elsewhere in UBT documents) sets
        the $\psi$-circle holonomy to the identity, making this
        topological term zero.
    \end{enumerate}
    Combining all three arguments: a CS term is forbidden by parity
    (argument 1+2) and vanishes by holonomy triviality (argument 3).
  \end{proof}
\end{proposition}

\begin{observation}[No CS implies $k_{\mathrm{CS}} = 0$]
  \label{obs:no_CS_kCS0}
  If $\mathcal{S}_{\mathrm{CS}} = 0$ (Proposition~\ref{prop:no_CS}),
  then $k_{\mathrm{CS}} = 0$ and the Kac-Moody level is determined
  entirely by the kinetic term $\mathcal{S}_{\mathrm{kin}}$.
  \textit{Why this is useful:} it eliminates the possibility that $k$
  is shifted by an arbitrary integer from a CS term; whatever $k$
  comes from the kinetic term is the true level.
\end{observation}

\subsection{Minimal level from the free kinetic action}

The UBT kinetic action on $S^1_\psi$ (for fixed 4D spacetime)
restricted to the $g \in \mathrm{SU}(2)$ sector is
(Eq.~\eqref{eq:S_kin_g}):
\begin{equation}
  \mathcal{S}_{\mathrm{kin}}^{g,\psi}
  = r^2 \int_0^{2\pi} \mathrm{d}\psi
    \Bigl(-\tfrac{1}{2}\,\mathrm{Tr}_{\mathrm{fund}}[J_\psi^2]\Bigr).
  \label{eq:S_kin_g_H2}
\end{equation}
This has the form of a WZW action at level $k = 2\pi r^2$
(Eq.~\eqref{eq:k_from_r2}).

The $\widehat{\mathrm{SU}(2)}_1$ WZW model is distinguished as the
\textbf{free-field point}: it is equivalent (via bosonisation) to
one free Dirac fermion on the circle, or equivalently to three free
real bosons compactified at the SU(2) self-dual radius.
The $\widehat{\mathrm{SU}(2)}_1$ model is the unique WZW model with
no relevant deformations and minimal operator content.

\noindent\textit{Minimality principle (H2):}
Given that:
(a)~no CS term is present (Proposition~\ref{prop:no_CS}),
(b)~the UBT kinetic action $\mathcal{S}_{\mathrm{kin}}$ is
    a \emph{free} (quadratic) action in $\Theta$, and
(c)~the $\mathrm{SU}(2)$ current algebra on $S^1_\psi$ must have
    an integer level $k \geq 1$,

the \emph{minimal consistent} Kac-Moody level is $k = 1$.
Any higher level $k \geq 2$ would require either a non-trivial CS
term (excluded by (a)) or a rescaled vacuum amplitude $r^2 = k/(2\pi)$
with $k \geq 2$ (not forced by the free kinetic action).
\textit{Why $k=1$ is minimal:} the free kinetic action with no CS
generates, at the quantum level, the minimal current algebra that is
consistent with the $\mathrm{SU}(2)$ symmetry.
The coupling $k \geq 2$ would correspond to a strongly-coupled theory
with additional degrees of freedom not present in the free kinetic action.

\begin{conjecture}[H2: $k = 1$ from minimality]
  \label{conj:k1_H2}
  The Kac-Moody level of $\widehat{\mathrm{SU}(2)}$ on the UBT
  $\psi$-circle is $k = 1$, motivated by:
  \begin{enumerate}
    \item Absence of CS term (parity symmetry of UBT, proved).
    \item Freeness of the kinetic action (no higher-derivative or
      interaction terms on $S^1_\psi$, proved).
    \item Minimality: $k = 1$ is the unique level consistent with a
      free biquaternion field and no CS shift.
  \end{enumerate}
  Status: \textbf{[MOTIVATED CONJECTURE]}.
  What is proved: points (1) and (2).
  What is conjectured: that (1)+(2) uniquely force $k = 1$ rather
  than being consistent with \emph{any} integer $k$.
\end{conjecture}

\subsection{Why $k=1$ is not uniquely proved}

The gap in Conjecture~\ref{conj:k1_H2} is the following:
A free quadratic action with no CS term is consistent with
\emph{any} $k \geq 1$, because $k$ enters only through the overall
normalisation $r^2 = k/(2\pi)$.
Without knowing $r^2 = \langle\mathrm{Sc}[\Theta^\dagger\Theta]\rangle_\mathrm{vac}$
from an independent source, $k$ is not fixed.
The minimality principle — while natural — is an external assumption,
not a consequence of $\mathcal{S}[\Theta]$ alone.

\begin{remark}[Non-renormalization of $k$]
  \label{rem:nonrenorm}
  The WZW level $k$ is protected against perturbative renormalization:
  once fixed at the classical level, quantum corrections cannot change it.
  This means: \emph{if} $k = 1$ at the classical level (by whatever
  argument), it remains $k = 1$ quantum mechanically.
  This remark does not remove the obstacle (we still need to justify
  $k = 1$ classically), but it confirms that a classical derivation is
  sufficient.
\end{remark}

\textbf{H2 result summary:}
\begin{itemize}
  \item CS term is absent (parity + holonomy arguments).
    \textbf{[PROVED]}.
  \item $k$ from kinetic term alone: $k = 2\pi r^2$.
    \textbf{[PROVED]}.
  \item $r^2$ not determined by $\mathcal{S}_{\mathrm{kin}}$ alone.
    \textbf{[GAP]}.
  \item $k = 1$ from minimality principle.
    \textbf{[MOTIVATED CONJECTURE — not uniquely proved]}.
\end{itemize}

%--------------------------------------------------------------------
\section{Approach H3: Dynkin Index / Representation Content}
\label{sec:H3}
%--------------------------------------------------------------------

\textbf{Status: [DEAD END for $k=1$] — see Observations~\ref{obs:H3_fund}
and~\ref{obs:H3_adj}.}

\subsection{Kac-Moody level from representation theory}

An alternative approach to determining $k$ uses the
\emph{embedding index} (Dynkin index) formula.
For a set of modes contributing to the $\widehat{\mathrm{SU}(2)}$
current algebra, each in representation $\mathbf{r}_i$ with Dynkin
index $T(\mathbf{r}_i)$:
\begin{equation}
  k = \sum_i n_i\, T(\mathbf{r}_i),
  \label{eq:k_Dynkin}
\end{equation}
where $n_i$ is the multiplicity of representation $\mathbf{r}_i$.
\textit{Why this formula:} the level $k$ of the resulting Kac-Moody
algebra equals the sum of Dynkin indices of all contributing modes,
by the Kac-Moody current algebra construction from free fields.
For $\mathrm{SU}(2)$:
\begin{align}
  T(\mathbf{fund}) &= \tfrac{1}{2} \quad
    \text{(fundamental, spin-$\tfrac{1}{2}$)}, \label{eq:T_fund} \\
  T(\mathbf{adj}) &= 2 \quad
    \text{(adjoint, spin-1)}, \label{eq:T_adj} \\
  T(\mathbf{j}) &= j(2j+1)(j+1)/3 \quad \text{(spin-$j$ rep).}
    \label{eq:T_j}
\end{align}

\subsection{H3a: Imaginary modes in the fundamental representation}

The UBT field has three imaginary $\mathbb{C}$-valued components:
$\Phi_k = \Theta_k$ for $k = 1, 2, 3$ (the $\mathbf{i},\mathbf{j},\mathbf{k}$
coefficients).
Under $\mathrm{SU}(2)$, these three modes transform in the
\emph{adjoint} (triplet) representation of $\mathrm{SU}(2)$, not the
fundamental.
However, examining the modes as complex scalars, one can ask whether
they can be combined into fundamental-representation doublets.

Suppose, for the purpose of this approach, that we treat each
$\Phi_k$ as contributing to the $\widehat{\mathrm{SU}(2)}$ current
algebra in the fundamental representation (perhaps as complex fermions
via bosonisation):
\begin{equation}
  k_{\mathrm{H3a}}
  = N_{\mathrm{phases}} \times T(\mathbf{fund})
  = 3 \times \tfrac{1}{2}
  = \tfrac{3}{2}.
  \label{eq:k_H3a}
\end{equation}
Computing the central charge:
\begin{equation}
  c_{\mathrm{SU}(2),3/2}
  = \frac{3 \times (3/2)}{(3/2) + 2}
  = \frac{9/2}{7/2}
  = \frac{9}{7}
  \approx 1.286.
  \label{eq:c_H3a}
\end{equation}
The predicted $B_{\mathrm{base}}$:
\begin{equation}
  B_{\mathrm{base}}^{\mathrm{H3a}}
  = c \times N_{\mathrm{eff}}^{3/2}
  = \tfrac{9}{7} \times 41.57
  \approx 53.4
  \;\neq\; 41.57.
  \label{eq:B_H3a}
\end{equation}

\begin{observation}[H3a: $k = 3/2$ fails]
  \label{obs:H3_fund}
  Treating the three imaginary modes $\Phi_{1,2,3}$ as fundamental-rep
  contributions gives $k = 3/2$, $c = 9/7 \approx 1.286$, and
  $B_{\mathrm{base}} \approx 53.4 \neq 41.57$.
  The half-integer level $k = 3/2$ is allowed for $\mathrm{SU}(2)$
  but does not reproduce the target.
  \textbf{[DEAD END for $k=1$]}.
\end{observation}

\subsection{H3b: Imaginary modes in the adjoint representation}

The natural representation of $\mathrm{Im}(\mathbb{H}) = \mathrm{span}\{\mathbf{i},\mathbf{j},\mathbf{k}\}$
under $\mathrm{SU}(2)$ is the \emph{adjoint} (triplet, spin-1):
\begin{equation}
  k_{\mathrm{H3b}}
  = N_{\mathrm{phases}} \times T(\mathbf{adj})
  = 3 \times 2
  = 6.
  \label{eq:k_H3b}
\end{equation}
Computing the central charge:
\begin{equation}
  c_{\mathrm{SU}(2),6}
  = \frac{3 \times 6}{6 + 2}
  = \frac{18}{8}
  = \frac{9}{4}
  = 2.25.
  \label{eq:c_H3b}
\end{equation}
\begin{equation}
  B_{\mathrm{base}}^{\mathrm{H3b}}
  = 2.25 \times 41.57
  \approx 93.5
  \;\neq\; 41.57.
  \label{eq:B_H3b}
\end{equation}

\begin{observation}[H3b: $k = 6$ fails]
  \label{obs:H3_adj}
  The adjoint-representation assignment (correct under the natural
  $\mathrm{SU}(2)$ action on $\mathrm{Im}(\mathbb{H})$) gives
  $k = 6$, $c = 9/4 = 2.25$, and $B_{\mathrm{base}} \approx 93.5 \neq 41.57$.
  \textbf{[DEAD END]}.
\end{observation}

\subsection{H3c: Can any representation assignment give $k=1$?}

For $k = 1$ from Eq.~\eqref{eq:k_Dynkin}:
$\sum_i n_i T(\mathbf{r}_i) = 1$.
The options with small integers are:
\begin{center}
\begin{tabular}{lllll}
  \toprule
  $n_i$ & Rep & $T(\mathbf{r}_i)$ & $k$ & Natural for UBT? \\
  \midrule
  $2$ & fund ($j=\tfrac{1}{2}$) & $\tfrac{1}{2}$ & $1$ & No: only 3 or 4 modes \\
  $1$ & spin-$j$ with $T=1$ & $T(j)=1$ requires $j\approx 0.77$ & $1$ & No: $j$ is non-integer \\
  $1$ & $\mathbf{adj}$ ($j=1$) & $2$ & $2$ & No \\
  \bottomrule
\end{tabular}
\end{center}
None of the natural representation assignments from $\mathbb{C}\otimes\mathbb{H}$
(3 imaginary modes in adjoint, or 4 complex components in bifundamental)
give $k = 1$ through Eq.~\eqref{eq:k_Dynkin}.
The only option with $n_i = 2$, $T = 1/2$ would require exactly
2~fundamental-rep modes, but UBT has 3~imaginary modes, not~2.

\textbf{H3 result summary:}
\begin{itemize}
  \item H3a (fundamental assignment):
    $k = 3/2$, $c = 9/7$, $B \approx 53.4 \neq 41.57$.
    \textbf{[DEAD END]}.
  \item H3b (adjoint assignment, natural):
    $k = 6$, $c = 9/4$, $B \approx 93.5 \neq 41.57$.
    \textbf{[DEAD END]}.
  \item H3c (search for $k=1$):
    no natural assignment from $\mathbb{C}\otimes\mathbb{H}$ gives $k=1$.
    \textbf{[DEAD END]}.
\end{itemize}

%--------------------------------------------------------------------
\section{Summary and Gap Status}
\label{sec:summary}
%--------------------------------------------------------------------

\begin{center}
\begin{tabular}{llll}
  \toprule
  Approach & Method & Conclusion & Status \\
  \midrule
  H1 (via $V_{\mathrm{eff}}$) & $k = 2\pi r^2$, $r^2 = f(n^*)$ & circular through $n^*=137$ & \textbf{[DEAD END (circular)]} \\
  H1 (canonical norm.)  & $r^2 = 1 \Rightarrow k = 2\pi \notin \mathbb{Z}$ & non-integer level & \textbf{[DEAD END]} \\
  H2 (absence of CS)    & CS absent (parity), $k_{\mathrm{CS}}=0$ & proved & \textbf{[PROVED]} \\
  H2 (minimality)       & free action + no CS $\Rightarrow k \geq 1$; min.\ $k=1$ & motivated, not proved & \textbf{[MOTIVATED CONJECTURE]} \\
  H3a (fundamental)     & $k = 3 \times \tfrac{1}{2} = \tfrac{3}{2}$; $B \approx 53.4$ & wrong & \textbf{[DEAD END]} \\
  H3b (adjoint)         & $k = 3 \times 2 = 6$; $B \approx 93.5$ & wrong & \textbf{[DEAD END]} \\
  H3c (search $k=1$)    & no natural assignment & no match & \textbf{[DEAD END]} \\
  \bottomrule
\end{tabular}
\end{center}

\subsection{What has been gained}

Comparing to the state in \texttt{b\_base\_g\_approaches.tex} (v65):
\begin{itemize}
  \item The free-boson / minimal-coupling motivation for $k=1$
    mentioned in v65 has been made more precise:
    \begin{enumerate}
      \item The formula $k = 2\pi r^2$ is now derived exactly from
        $\mathcal{S}[\Theta]$ and the biquaternion scalar-part identity.
      \item The absence of a CS term is now argued from parity and
        holonomy, not just asserted.
      \item The WZW non-renormalization theorem confirms that $k=1$, once
        established, cannot be shifted by loops.
    \end{enumerate}
  \item The obstacle is now precisely identified: $k=1$ requires
    $r^2 = 1/(2\pi)$, which is not determined by
    $\mathcal{S}_{\mathrm{kin}}$ alone and cannot be obtained without
    either invoking $n^*$ (circularity) or the minimality principle
    (extra assumption).
  \item The representation-theory route (H3) is exhausted and declared
    a dead end.
\end{itemize}

\subsection{Updated gap status}

\begin{gap}[Gap (G3-k): Level of $\widehat{\mathrm{SU}(2)}$ in UBT — Updated]
  \label{gap:G3_k_updated}
  The formula $k = 2\pi\,\langle\mathrm{Sc}[\Theta^\dagger\Theta]\rangle_\mathrm{vac}$
  has been derived (Eq.~\ref{eq:k_from_r2}).
  The CS term is proved absent (Proposition~\ref{prop:no_CS}).
  The level $k = 1$ is strongly motivated by the minimality principle
  (absence of CS, free action, $k=1$ as the unique level of the free-field
  point of the $\widehat{\mathrm{SU}(2)}$ WZW model).
  However, a proof that $\langle\mathrm{Sc}[\Theta^\dagger\Theta]\rangle_\mathrm{vac}
  = 1/(2\pi)$ from $\mathcal{S}[\Theta]$ without circular use of $n^*$
  remains open.
  Status: \textbf{[OPEN — $k=1$ is [MOTIVATED CONJECTURE],
  stronger than [PARTIAL] in v65, but not yet [PROVED]]}.
\end{gap}

\subsection{If $k=1$ is accepted as a motivated conjecture}

Under Conjecture~\ref{conj:k1_H2}, $c_{\widehat{\mathrm{SU}(2)},1} = 1$, and:
\begin{equation}
  B_{\mathrm{base}}
  = c_{\widehat{\mathrm{SU}(2)},1} \times N_{\mathrm{eff}}^{3/2}
  = 1 \times 12^{3/2}
  = 41.569\ldots
  \label{eq:B_base_final}
\end{equation}

\begin{theorem}[$B_{\mathrm{base}} = N_{\mathrm{eff}}^{3/2}$ — Status Upgrade]
  \label{thm:B_base}
  \textbf{Status: [MOTIVATED CONJECTURE → PARTIAL (H2)]}.
  Under the motivated conjecture that the UBT $\psi$-circle carries
  a $\widehat{\mathrm{SU}(2)}_1$ Kac-Moody symmetry ($k=1$, motivated by
  parity, free kinetic action, and minimality via H2), the central charge
  is $c = 1$ and:
  \[
    B_{\mathrm{base}}
    = c \cdot N_{\mathrm{eff}}^{3/2}
    = N_{\mathrm{eff}}^{3/2}
    \approx 41.57.
  \]
  Inputs: $N_{\mathrm{eff}} = 12$ (proved), $k=1$ (motivated conjecture).
  Output: $B_{\mathrm{base}} = 41.57$ (no free parameters given $k$).
  The formula $B_{\mathrm{base}} = N_{\mathrm{eff}}^{3/2}$ is consistent with
  the Gaussian path-integral argument (Approach~A2) and now also with
  the Kac-Moody argument (Approach~G3/H2).
  This double consistency strengthens, but does not prove, the result.
\end{theorem}

%--------------------------------------------------------------------
\section{In Plain Language}
\label{sec:plain}
%--------------------------------------------------------------------

\textbf{In plain language:}
We investigated whether the integer $k$ characterising the quantum symmetry
of the UBT imaginary-time circle (called the Kac-Moody level) can be
computed directly from the UBT action, with the goal of confirming $k = 1$.

The first approach (H1) expresses $k$ as proportional to the squared amplitude
of the UBT field in the vacuum.
The formula is clean ($k = 2\pi r^2$), but to use it we need to know $r^2$,
which is set by the minimum of the effective potential — and that minimum is
at $n^* = 137$, the very quantity we are trying to derive.
Trying to use the natural normalization $r = 1$ instead gives $k = 2\pi \approx 6.28$,
which is not a valid integer level.
This approach hits a wall in both directions.

The second approach (H2) asks: is the action allowed to contain a
Chern-Simons term, which could shift $k$ to any integer?
We show that it cannot: a Chern-Simons term would break parity, but UBT
has parity symmetry (the two vacuum sectors at $n^* = 137$ and $n^{**} = 139$
are mirror images of each other).
Independently, the gauge holonomy around the imaginary-time circle is trivial,
making any would-be topological Chern-Simons contribution vanish.
With no Chern-Simons term, the level is set entirely by the kinetic part of
the action.
A free kinetic action with no interactions and no Chern-Simons term is
naturally at the lowest possible level, $k = 1$.
This is a physical principle — the simplest quantum algebra consistent with
the symmetry — not a theorem, so it remains a motivated conjecture.

The third approach (H3) tries to count how many modes contribute to the
symmetry algebra, and multiply by their algebraic weight.
The three imaginary components of the UBT field, transforming as a triplet
under the symmetry, naturally give $k = 6$ (adjoint weight 2, three modes) or
$k = 3/2$ (fundamental weight $1/2$, three modes).
Neither gives $k = 1$, and no natural assignment from the biquaternion algebra
produces $k = 1$.
This approach rules out $k = 1$ from representation counting alone.

After all three approaches, $k = 1$ is best described as a
\emph{motivated conjecture}: it is the simplest consistent choice for a free
field without Chern-Simons terms, it is algebraically clean (giving central
charge exactly $1$), and it is now confirmed to not be contradicted by the
representation-theory approach.
However, it has not been derived from the action as a theorem.
The gap (G3-k) remains open, with the concrete next step being to compute the
vacuum amplitude $r^2 = \langle\mathrm{Sc}[\Theta^\dagger\Theta]\rangle_\mathrm{vac}$
from the effective potential without using $n^*$ as input.

\ifdefined\INCLUDEMODE\else
\end{document}
\fi
